% Generated draft: requires mathematical and source review.
\documentclass[11pt]{article}
\usepackage[margin=25mm]{geometry}
\usepackage{fontspec}
\setmainfont{Latin Modern Roman}
\usepackage{amsmath,amssymb,mathtools,unicode-math}
\setmathfont{Latin Modern Math}
\usepackage{enumitem,needspace,xurl,hyperref}
\setlength{\parindent}{0pt}
\setlength{\parskip}{5pt}
\setlength{\emergencystretch}{3em}
\newcommand{\N}{\mathbb{N}}
\newcommand{\Z}{\mathbb{Z}}
\newcommand{\Q}{\mathbb{Q}}
\newcommand{\R}{\mathbb{R}}
\newcommand{\C}{\mathbb{C}}
\newcommand{\F}{\mathbb{F}}
\newcommand{\A}{\mathbb{A}}
\newcommand{\PP}{\mathbb{P}}
\newcommand{\E}{\mathbb{E}}
\newcommand{\eps}{\varepsilon}
\newcommand{\dd}{\,\mathrm{d}}
\begin{document}
\section*{UnsolvedMath ID: 1900003}
\textbf{Source problem code:} catalogue entry
\subsection{Definitions and mathematical statement}
The title and research summary identify the objects in this question as faces of sails in the geometry of multidimensional continued fractions. The literal statement is: ``Classify all combinatorial possible types of faces.'' It supplies no definition of the sails or admissible cones.

For context, a customary sail construction starts with a full-rank lattice $L\subset\R^d$ and a pointed convex cone $C$ with vertex $0$, takes the convex hull of nonzero lattice points in the chosen cone region, and considers its boundary facing the origin. For example, once closed-cone and closure conventions are fixed, one can write
\[
 K(C,L)=\overline{\operatorname{conv}}\bigl((C\cap L)\setminus\{0\}\bigr).
\]
The sail consists of the appropriate boundary faces of this hull. Conditions such as simpliciality, irrationality of boundary hyperplanes, and whether cone-boundary lattice points are included are part of the missing source setup; this contextual construction does not select them on the source's behalf.

Two polyhedral faces have the same \emph{combinatorial type} when their face posets, ordered by inclusion, are isomorphic. This records face incidences rather than metric shape or lattice equivalence. If $\mathscr S_d$ denotes the intended class of $d$-dimensional sail constructions and $\mathscr F(S)$ the intended faces of a sail $S$, the classification target can be expressed as
\[
 \left\{[\mathcal P(F)]_{\cong}:d\text{ in the intended dimension range},\ 
 S\in\mathscr S_d,\ F\in\mathscr F(S)\right\},
\]
where $\mathcal P(F)$ is the face poset of $F$. A complete classification would characterize exactly which types occur and account for their realizability by the allowed sails. The supplied summary describes the problem as one in arbitrary dimension, but the source excerpt does not determine $\mathscr S_d$, whether only facets or all-dimensional faces are intended, or whether all relevant faces are bounded. These gaps prevent a unique fully quantified classification problem.

The registry records OPEN-TRIAGE with moderate confidence: its summary says that no complete classification was found and that low-dimensional partial results exist, without identifying those results. This is supplied status evidence rather than independent verification.
\subsection{Short English statement}
Classify the face-incidence types that can occur among faces of multidimensional continued-fraction sails. The excerpt omits the underlying sail construction and the intended face and dimension ranges, so those parameters must be recovered before a precise exhaustive classification can be stated.
\subsection{Sources}
\begin{itemize}
\item Local UnsolvedMath registry, permanent ID 1900003, source code catalogue entry, ``Geometry of Continued Fractions --- Faces of sails''; supplied status open, classification OPEN-TRIAGE, updated 2026-08-06. The research summary qualifies its assessment as moderate confidence and mentions unspecified low-dimensional partial results.
\item Karpenkov, \emph{Open problems in geometry of continued fractions} (2017), Problem 3, \url{https://arxiv.org/abs/1712.01450}. The record reports source-TeX extraction and access on 2026-07-29. It says the source presents an open item and current status requires release review. Rights: NEEDS\_REVIEW; public arXiv source TeX inspected according to the record.
\end{itemize}
\end{document}
